\capitulo{5}{What Does the Classifier Need to Remember?}
I hand you a ticket and a pencil. The ticket came from one of two printers, A or B, chosen at random with equal chances. It has two windows covered by paper flaps. In each window the machine printed heads or tails. Your job is to identify the machine.

Open the first flap. Heads. Half of A's tickets show heads in the first window. For B it's half as well. What's your bet on A? Write $1/2$ on the back. Now do me a favor. Since that head did nothing to change your bet, take your pencil and cross out the picture until you can't see it anymore.

If you're still holding the pencil, wait a moment.

There's a difference between the printers. A chooses the first face at random and, nine times out of ten, copies it into the second window. B also chooses the first at random, but nine times out of ten it prints the opposite face in the second window. Those chances of copying are the same whether the first window shows heads or tails. Someone who sees only the second window will also find the same frequency of heads in both machines. Under A, for example, the first could have been a head and the machine copied it, or the first could have been a tail and the machine reversed it. Add them up and we have $(1/2)(9/10)+(1/2)(1/10)=1/2$. Under B the two copying and reversing probabilities simply trade places.

Close the first flap without crossing out the picture. Before you open the other flap, write down two bets on A. If I show you another head, you'll have seen a repeat. The odds in that case are nine to one in favor of A. Your bet on A will be $90\%$. If I show you a tail, it will be $10\%$. Now imagine the first picture had been tails. The back of your ticket would still say $1/2$, but you'd have to reverse the instructions for the second window.

Open the second flap. Heads. Someone who just walked in and saw only that window gives A a probability of $1/2$. You give it $9/10$. When did you learn something from the first head? You could already show what it had taught you \emph{before} opening the second flap. It was there in the two bets you just wrote down.

\secao{A model that counts every head and misses the repeats}
Let's give the problem to two learners. Neither knows the printing trick. Each gets a notebook with two hundred tickets whose origins we know. I've chosen counts that reproduce the proportions in the workshop exactly. As before, $H$ stands for heads and $T$ for tails, this time in the order of the two windows.
\begin{center}
\begin{tabular}{lrrrrr}
\toprule
Machine & $HH$ & $HT$ & $TH$ & $TT$ & Total \\
\midrule
A & 45 & 5 & 5 & 45 & 100 \\
B & 5 & 45 & 45 & 5 & 100 \\
\bottomrule
\end{tabular}
\end{center}

The first learner wants to make tickets like the ones it found. It measures the frequency of heads in each window for each machine. It gets $1/2$ four times. If I order a hundred A tickets, it will draw the two faces independently at those rates. It expects to produce twenty-five $HH$ tickets. A produced forty-five. If I order a hundred B tickets, it expects the same twenty-five. B produced five. Which of the four counts did it get wrong?

It\bookterms{Naive Bayes|Conditional independence} got them all right. The learner itself put the instruction to draw the faces independently into the program. When we use a generator like this for classification, we call it \emph{Naive Bayes}. Here is the assumption we put in by hand,
\[
     p(x_1,x_2\mid y)=p(x_1\mid y)\,p(x_2\mid y).
\]
Why would anyone agree to it? Add one more window to every pattern you already have. Each old pattern now has two possible continuations. Do that doubling for ten windows and we arrive at $2^{10}=1{,}024$ patterns. If I give you the probabilities of 1,023 of them, you get the last one from what's left over to make one. That's $1{,}023$ free numbers per machine. Our learner asks for ten rates of heads per machine, one for each window, and it gets to use every ticket to estimate every rate. If I know that each print head has its own independent randomizer, and that its settings stay fixed within each machine, I have a reason to try that shortcut.

Add the two rows of the table. There are fifty tickets in each of the four columns. Anyone who checks the whole notebook for independence will find it. But there's a vertical bar in the formula, with a $y$ on its right. Stay in row A. After a head, the second window shows heads in $45/50$ of the cases. After a tail, in $5/50$. The assumption the program needed fails precisely when we check it in the group it describes.

We can write a different recipe for the generator. First draw the face on the left. Then draw whether the machine will copy it. To train that second draw, just add the two outer columns. Ninety of A's hundred tickets repeat the first face, compared with ten of B's hundred. On an order of a hundred A tickets, the new recipe expects $100(1/2)(9/10)=45$ pairs marked $HH$. When it receives an unlabeled $HH$ ticket, it assigns A a probability of $45/(45+5)=9/10$. We still have a generative model. What we changed was an assumption about the printer.

\secao{What a multiplication puts into a model}
The second learner wants to estimate the probability of the letter directly from the faces. It writes $+1$ for heads and $-1$ for tails. It proposes a logistic regression in which the log odds of A are $b_0+b_1x_1+b_2x_2$. Let's try its formula without fitting a single coefficient.

On an $HT$ ticket, change the tail on the right to a head. The bet on A should go up from $1/10$ to $9/10$. Make the very same change on a $TT$ ticket. This time the bet should go down from $9/10$ to $1/10$. Cover the left window in both calculations. Changing the right from $-1$ to $+1$ takes $-b_2$ out of the sum and puts $+b_2$ in its place. The change in log odds is $2b_2$ on both tickets. Which sign would you choose for $b_2$?

The learner can keep collecting tickets. It will still need a sign that changes depending on the first window. I want an instruction that leaves the reading on the right alone when the left shows $+1$ and reverses its sign when the left shows $-1$. Multiplying by the reading on the left does just that. The product $z=x_1x_2$ gives $+1$ for repeats and $-1$ for mismatches. The odds in favor of A in those two groups are nine to one and one to nine. On the logarithmic ruler those marks lie at the same distance from zero, one on either side. Use $z$ to choose the side and $\ln 9$ to mark the distance. Fitting a logistic regression to these proportions gives
\[
              \ln\frac{p(A\mid z)}{p(B\mid z)}=z\ln 9 .
\]
The product is an \emph{interaction term}. Multiply $x_2$ by $x_1$ and watch what happens to the contribution of a head on the right when you change the face on the left. That's the change of sign we asked for. Even this fix involved a choice of our own. By supplying only $z$, we lumped $HH$ together with $TT$. The program lost the freedom to give them different probabilities. So I'd keep the original faces in the next batch to check whether the distinction I erased can still safely be ignored.

The generator and the logistic model agree on the letter for \emph{every} ticket. Can I put away the generator and use logistic regression to plan production? On both machines, change the rate of heads in the first window from $1/2$ to $0.8$, keeping the copying probabilities the same. For every hundred A tickets we now expect 72 pairs marked $HH$. For every hundred B tickets we expect eight. The classifier still gives A a probability of $72/(72+8)=9/10$ on seeing $HH$. For the other patterns, too, the changed frequency of the first face cancels when we take the ratio between the machines. Anyone ordering a hundred A tickets has good reason to ask whether to expect 45 or 72 pairs of heads. The logistic model can price every bet on the letter correctly without discovering what I changed on the printer.

\secao{Who chose the examples?}
I often tell my students that we do the modeling when we decide what assumptions we have to bring to the problem ourselves. The next two workshops want the same thing. Tomorrow an unlabeled ticket will arrive and someone will have to guess its origin. In which workshop would I start with each approach?

In\bookterms{Sampling|Discriminative learning} the first, it costs money to discover the letter. The workshop investigates only half the tickets with matching faces and one in ten of those with different faces. Within each group it chooses the tickets at random. Apply that procedure to our table. Of the hundred matching tickets, we expect to investigate fifty, of which 45 will be A's. Of the hundred mismatches we expect to investigate ten, of which one will be an A. The expected audit counts put the share of matching tickets at $50/60$. On the conveyor belt it's one half. If I used this file to generate another print run, I'd already be wrong on my first order. But the chance of A among the matches is still $45/50$, and among the mismatches it's $1/10$. Sampling shrank each column without favoring either letter within it. I'd start with the discriminative approach because I understand why the proportions I need survived sampling. The decisive information came from whoever chose the tickets for inspection.

In\bookterms{Generative learning} the other workshop I go to the two printers myself. I take a hundred tickets at random from A's output and a hundred from B's. Now I can estimate what each machine prints. My notebook is half A because I just decided it would be. On tomorrow's shift only one ticket in ten will come from A. If the mechanic tells me the settings are staying put, I can carry over the rates I measured at the bench and use the production mix. A repeat still favors A by a factor of nine. Since the starting odds have fallen to one to nine, after a repeat my bet on A is $1/2$. Here I'd start with the generator. Stable settings are an assumption I'd want to be able to defend with the mechanic. The hundred tickets from each machine can't make that promise on their own.

You can also correct the bet of a calibrated discriminative model if you know the old mix, the new mix, and that the same stability holds. In practice, the choice depends on what I can justify about the next shift. To test my choice, I'd set aside a batch from the actual conveyor belt. If the job is to sort tickets, I'd uncover the labels to measure the cost of sending a ticket to the wrong bin. If it's to plan the print run, I'd compare the number of tickets of each pattern with the quantities I'd promised the buyer.

\begin{quadro}{Before you erase a column}
\itemq{What I need to find out} When sampling preserves the share of each letter within each pattern and my task is to identify the machine, I can fit that share directly. When I know each machine's output and the mix for the next shift, I can put the prediction together from the pieces I measured.
\itemq{What I put in by hand} Drawing the two faces separately and representing them only by their product are different assumptions. One rules out dependence within each machine. The other asks for the same bet on two heads as on two tails.
\end{quadro}

Back to your ticket. Can I build a reader with room to store only one bit? As it opens the first flap it writes heads or tails in that space. When it opens the second, it compares what it sees with what it stored, then replaces the first entry with the answer to the question \emph{same or different?} Now it can cross out both faces and still bet on the machine. We had enough room from the beginning. What had to change between the two readings was the meaning of that one entry.

\begin{quadro}{What to take away from this chapter}
\itemq{1} A face can be needed to interpret the next one, even when it leaves our current bet on the machine unchanged.
\itemq{2} The independence in Naive Bayes is conditional on the class. Multiplying the two faces captures the example's interaction, but also groups patterns that the model must then treat alike.
\itemq{3} Data collection determines which proportions we can carry into the next shift. Making that choice requires knowing the workshop and checking predictions on fresh data.
\itemq{What changed} We now assess a representation by what it lets us answer and what it makes us forget. Even a single bit can serve, provided its meaning keeps up with the reading.
\itemq{To carry forward} Have two classifiers that always announce the same letter learned the same thing?
\end{quadro}
